\clearpage
\begin{center}
    {\Large \textbf{ABSTRACT}\par}
\end{center}
\vspace{1.0cm}

Continuous Integration (CI) is a foundational software engineering practice wherein developers merge code changes frequently. However, running comprehensive regression test suites on every pull request is prohibitively time-consuming and computationally expensive. While Machine Learning-based Regression Test Selection (ML-RTS) offers a promising heuristic approach to predict test failure likelihoods, existing models exhibit systematic overconfidence, lack uncertainty awareness on out-of-distribution code refactorings, and risk silent regression bug escapes into production.

In this project, we propose and implement \textbf{ConfTest}, a confidence-calibrated, uncertainty-aware selective regression test selection framework. ConfTest mines a canonical 32-dimensional feature space spanning AST syntactic complexity, static NetworkX call-graph dependency coupling, code churn metrics, and historical test failure telemetry. We construct a 5-seed deep ensemble to quantify epistemic uncertainty ($\sigma$) and apply post-hoc Temperature Scaling to calibrate predicted failure probabilities, reducing Expected Calibration Error from $0.0390$ to $0.0242$ at $T^* = 0.7941$ (paired improvement $0.0148$, $95\%$ CI $[0.0030, 0.0393]$). Isotonic regression was also fitted and rejected: it attains a lower mean ECE but nearly doubles worst-bin error, and an abstention gate is thresholded on the tail.

Furthermore, ConfTest introduces an economic risk-coverage selective prediction policy that executes targeted fast test subsets when confidence is high, while autonomously triggering full-suite fallback when epistemic divergence exceeds threshold bounds. Our empirical findings are mixed and we report them as such. On $183$ held-out commits ($86{,}469$ commit--test pairs, all labels obtained by executing the suite against generated mutants) the shipped operating point observes \textbf{$100.0\%$ regression failure recall with zero escaped commits}, but abstains on $97.81\%$ of commits and therefore reduces test executions by only $\mathbf{3.11\%}$, with a $\mathbf{0.0\%}$ wall-clock reduction. This is an empirical result, not a formal guarantee for future commits or unseen repositories. Tightening the abstention threshold to reach a $32.85\%$ reduction drops held-out recall to $87.69\%$ with $15$ escaped commits, failing the $95\%$ recall floor this project pre-registered by an $8.33$ percentage-point validation-to-test generalization gap. We conclude that the selective-prediction mechanism is reusable, while the underlying ranker (PR-AUC $0.1393$ against a $3.71\%$ positive rate) is not yet accurate enough to convert abstention into a useful saving. A research-prototype stack comprising a FastAPI backend, GitHub webhook PR bot, and Streamlit visual analytics dashboard was implemented.

\vspace{0.5cm}
\noindent \textbf{Keywords:} Regression Test Selection, Confidence Calibration, Uncertainty Estimation, Continuous Integration, Temperature Scaling, Deep Ensembles.

\clearpage
